\section{\sn: Hierarchical Planning and Task-Specific Agents}
\label{sec:arch}

% \subsection{Deficiencies in Single Agents}

As mentioned, ReAct-style agents iterate by taking actions, observing the
response, and repeating. Although successful for many kinds of tasks, the
repeated iteration can make long-term planning for cybersecurity tasks
fail because 1) the context can extend rapidly for cybersecurity tasks, and 2) 
it can be difficult for the LLM to try many different exploits. For example, 
prior work has shown that if an LLM agent attempts one type of vulnerability, 
backtracking to try another type of vulnerability is challenging for a single 
agent \cite{fang2024llm2}.

One method of improving the performance of a single agent is to use multiple
agents. In this work, we introduce a method of using hierarchical planning and
task-specific agents (\sn) to perform complex, real-world tasks.

\subsection{Overall Architecture}

\begin{figure}
  \includegraphics[width=\columnwidth]{figures/pdfs/llm-0d-hacking}

  \caption{Overall architecture diagram of \sn. We have other task-specific,
  expert agents beyond the ones in the diagram.}
  \label{fig:arch}
\end{figure}

\sn has three major components: a hierarchical planner, a set of task-specific,
expert agents, and a team manager for the task-specific agents. We show an
overall architecture diagram in Figure~\ref{fig:arch}.

Our first component is the hierarchical planner, which explores the environment
(i.e., websites). After exploring the environment, it determines the 
set of instructions to send to the team manager. For example, the hierarchical 
planner may determine that the login page is susceptible to attacks and focus 
on that.

Our second component is a team manager for the task-specific agents. It
determines which specific agents to use. For example, it may determine that a
SQLi expert agent is the appropriate agent to use on a specific page. Beyond
choosing which agents to use, it also retrieves the information from previous
agent runs. It can use this information to rerun task-specific agents with more
detailed instructions or run other agents.

Finally, our last component is a set of task-specific, expert agents. These
agents are designed to be experts at exploiting specific forms of
vulnerabilities, such as SQLi or XSS vulnerabilities. We describe the design of
these agents below.


\subsection{Task-Specific Agents}

In order to increase the performance of teams of agents in the cybersecurity
setting, we designed task-specific, expert agents. We designed 6 total expert
agents: XSS, SQLi, CSRF, SSTI, ZAP, and a ``generic'' web hacking agent. Our AI
agents have: 1) access to tools, 2) access to documents, and 3) specific prompts.

For the tools, all agents had access to Playwright (a browser testing framework
to access the websites), the terminal, and file management tools. The ZAP agent
also had access to ZAP \cite{bennetts2013owasp}, while the SQLi agent 
had access to sqlmap \cite{sqlmap}. The agents accessed the
websites via Playwright. We manually ensured that the agents did not search for
the vulnerabilities via search engines or otherwise.

To choose the documents, we manually scraped the web for relevant documents for
the specific vulnerability at hand. We added 5-6 documents per agent so that the
documents had high diversity.

Finally, for the prompt, we used the same prompt template. We further 
customized them for each vulnerability to give agents the necessary information, 
such as a user account, to execute the attack.

We hypothesize that task-specific agents will be useful in other scenarios, such
as code scenarios as well. However, such an investigation is outside the scope
of this work.


\subsection{Implementation}

In our specific implementation for \sn for web vulnerabilities, we used the
LangChain and LangGraph library in conjunction to APIs of Fireworks and OpenAI 
assistants. We used LangGraph's functionality to create a graph of agents and 
passed messages between agents using LangGraph. The individual agents were 
implemented with a conjunction of OpenAI Assistants, Fireworks, and LangChain.

To reduce the token count (directly reducing costs), we observed that the
client-side HTML was the vast majority of the tokens. We implemented an HTML
simplifying strategy to reduce this cost. Before passing the HTML of the webpage
to the agent, we remove unnecessary HTML tags (such as image, svg, style, etc.)
tags that are irrelevant to the agent.
